\chapter{Судалгааны арга зүй}
\label{ch:method}

%======================================================================
% ЭНЭ ФАЙЛД БАЙГАА БҮХ ЗҮЙЛ БАГШИД ҮЗҮҮЛЭХ НООРОГТ ХЭВЛЭГДЭНЭ (бүтэн ноорогт ч мөн): 3.1–3.5.
% Зөвхөн бүтэн ноорогт орох үнэлгээ, яриа таних аргын сонголт:
% Chapters/full/ch3_evaluation_and_stt_choice.tex (файлын төгсгөлийн \fullinput мөр).
%
% 2026.11.16 (3), хэрэглэгчийн шүүмжээр: (1) мэдлэгийн мод бол энэ ажлын RAG-ийн суурь — 3.4-т
% бүтэц, зураг, хайлтад хэрхэн ашиглахыг алхмаар тайлбарласан; (2) «ирээгүй, байхгүй, хоосон» гэж
% бичихгүй — цуглуулах мэдээллийг «Цаашдын ажил»-д «цуглуулна» гэж бичнэ; (3) «сургуулийн албад»
% биш «сургуулийн зүгээс өгсөн мэдээлэл», «сургуулийн ажилтнууд».
% Кодтой тулгасан (11/16): config/knowledge-topics.json (үндэс, 11 хүрээ, 31 дэд сэдэв, хаягаар оноох
% дүрэм), config/knowledge-domains.json (түлхүүр үг, enabled=false), website_retrieval.py
% (predict_domain, тэнцвэл оноо нэмэхгүй), config/referral.json (topic_offices approved=false).
% Өмнөх хувилбарууд: thesis/archive-20261116/.
%======================================================================

\section{Судалгааны арга барил}
\label{sec:approach}

Судалгаа дөрвөн ажлаас бүрдэнэ. Нэгдүгээрт, сургуулийн ажилтнуудтай уулзаж, системд тавигдах шаардлагыг тодорхойлсон. Хоёрдугаарт, сургуулийн мэдээллийг вэб сайтаас болон сургуулиас авч, мэдлэгийн мод хэлбэрээр зохион байгуулсан. Гуравдугаарт, мэдлэгийн модонд суурилсан хайлтаар асуултад хариулах системийг боловсруулж хэрэгжүүлсэн. Дөрөвдүгээрт, системийг ресепшнд асуудаг бодит асуултаар үнэлнэ. Эхний гурван ажил зэрэгцэн явсан бөгөөд дөрөв дэх ажил бэлтгэлийн шатанд байна. Үнэлгээний хэмжүүрийг үр дүнг харахаас өмнө тогтооно.

\section{Системд тавигдах шаардлага}
\label{sec:requirements}

Ресепшний ажилтны мэдээлснээр (\tref{tab:reception_topics}) хэрэглэгч хоёр бүлэгт хуваагдана. Өдөр тутам ирдэг хүмүүс (оюутан, эцэг эх, ажлын байр сонирхогч) голчлон хэнтэй холбогдох, хаана очихыг асууна. Нээлттэй өдөрлөгөөр ирдэг хүмүүс (элсэгч, эцэг эх) сургалтын орчин, төлбөр, дотуур байр зэрэг сургуулийг танилцуулах мэдээллийг асууна. Системийн хэрэглээний тохиолдлыг \fref{fig:usecase}-д харуулав.

\begin{figure}[ht]
\centering
\adjustbox{max width=\textwidth}{% Хэрэглээний тохиолдлын (use case) диаграмм (TikZ).
% Эх: demo_server.py-ийн нээлттэй API (/api/route, /api/voice/transcribe, /api/voice/speak,
% /api/admin/*) ба tools/import_school_tables.py. Оролцогчдыг 2026.11.12-ны ресепшний
% мэдээллээр хоёр бүлэгт хуваасан: өдөр тутмын зочин, нээлттэй өдөрлөгийн зочин.
% Байрлалыг (x,y) координатаар өгсөн.
\begin{tikzpicture}[
    font=\footnotesize,
    uc/.style={draw, ellipse, align=center, text width=3.3cm, minimum height=0.95cm, inner sep=1pt},
    actor/.style={align=center, font=\footnotesize},
    extbox/.style={draw, align=center, text width=2.5cm, minimum height=0.9cm, fill=black!8},
    line/.style={semithick},
    dep/.style={dashed, -{Stealth[length=2mm]}},
    gen/.style={-{Triangle[open, length=2.6mm, width=2.6mm]}, semithick},
    lab/.style={font=\scriptsize, inner sep=1.5pt, fill=white}]

% --- хүний дүрс зурах макро: #1 нэр, #2 (x,y), #3 бичиг
\newcommand{\stickman}[3]{%
  \begin{scope}[shift={(#2)}]
    \draw[semithick] (0,0.55) circle (0.17);
    \draw[semithick] (0,0.38) -- (0,-0.1) (-0.3,0.2) -- (0.3,0.2) (0,-0.1) -- (-0.25,-0.5) (0,-0.1) -- (0.25,-0.5);
    \node[actor, anchor=north] at (0,-0.55) {#3};
    \coordinate (#1) at (0.35,0.1);
    \coordinate (#1-l) at (-0.35,0.1);
  \end{scope}}

% --- системийн хүрээ
\draw[semithick] (2.4,2.1) rectangle (11.0,-9.6);
\node[font=\footnotesize\bfseries] at (6.7,1.75) {Сургуулийн ухаалаг туслах};

% --- хэрэглээний тохиолдлууд
\node[uc] (typed)  at (4.6,0.7)   {Бичгээр асуух};
\node[uc] (voice)  at (4.6,-0.8)  {Дуугаар асуух};
\node[uc] (answer) at (4.6,-2.5)  {Эх сурвалж, огноотой\\хариулт авах};
\node[uc] (refer)  at (4.6,-4.2)  {Хэнд, хаана хандахыг\\мэдэх};
\node[uc] (listen) at (4.6,-5.9)  {Хариултыг дуугаар\\сонсох};
\node[uc] (stt)    at (8.9,-0.8)  {Яриаг текст болгох};
\node[uc] (tts)    at (8.9,-5.9)  {Текстийг дуу болгох};
\node[uc] (import) at (8.9,-7.4)  {Сургуулийн хүснэгт\\оруулах};
\node[uc] (publish) at (8.9,-8.8) {Мэдээллийн шинэ хувилбар\\шалгаж нийтлэх};

% --- оролцогчид
\stickman{guest}{0.4,-2.4}{Зочин}
\stickman{daily}{-3.2,-0.4}{Өдөр тутмын зочин\\{\scriptsize оюутан, эцэг эх,}\\{\scriptsize ажил горилогч}}
\stickman{open}{-3.2,-4.4}{Нээлттэй өдөрлөгийн\\зочин\\{\scriptsize элсэгч, эцэг эх}}
\stickman{staff}{13.4,-8.0}{Сургуулийн\\ажилтан}
\node[extbox] (gemini) at (13.4,-0.8) {Яриа таних\\үйлчилгээ\\{\scriptsize Gemini 3.5 Transcribe}};
\node[extbox] (edge)   at (13.4,-5.9) {Дуу үүсгэх\\үйлчилгээ\\{\scriptsize Edge TTS}};

% --- холбоос
\foreach \u in {typed, voice, answer, refer, listen} \draw[line] (guest) -- (\u.west);
\draw[gen] (daily) -- (guest-l);
\draw[gen] (open) -- (guest-l);
\draw[dep] (voice) -- node[lab] {«include»} (stt);
\draw[dep] (listen) -- node[lab] {«include»} (tts);
\draw[dep] (refer) -- node[lab, right] {«extend»} (answer);
\draw[line] (stt) -- (gemini);
\draw[line] (tts) -- (edge);
\draw[line] (staff-l) -- (import.east);
\draw[line] (staff-l) -- (publish.east);
\end{tikzpicture}
}
\caption{Системийн хэрэглээний тохиолдлын диаграмм (ажилтны үйлдэл төлөвлөгдсөн)}
\label{fig:usecase}
\end{figure}

Сургалтын албатай хийсэн уулзалтаас дараах гурван шаардлагыг тодорхойлов.
\begin{enumerate}
    \item \textbf{Албан эх сурвалжтай хариулт.} Хариулт бүр сургуулийн албан мэдээллээс гарна. Хариулттай хамт эх сурвалж, мэдээллийн огноог харуулна.
    \item \textbf{Хэнд, хаана хандахыг заах.} Хариулт олдоогүй үед систем «олдсонгүй» гээд зогсохгүй. Тухайн мэдээллийг мэдэх алба эсвэл ажилтан, түүний байрлах өрөө, давхрыг зааж өгнө.
    \item \textbf{Нэг мэдээлэл.} Сургуулийн вэб сайт, Facebook хуудасны чатбот, энэ систем гурав ижил асуултад ижил хариулт өгнө.
\end{enumerate}

Эхний ба гуравдугаар шаардлагад нийцүүлэн систем хариултыг шинээр зохиохгүй, хайлтаар олдсон эх сурвалжийн өгүүлбэрийг өөрчлөлгүй харуулна. Хэлний загвараар хариулт бичүүлбэл ижил асуултад өгүүлбэр, заримдаа агуулга нь өөр хариулт гарч болно. Гуравдугаар шаардлагыг бүрэн хангахын тулд сайт, чатбот, систем гурав нэг баталгаажсан эх сурвалжаас унших шаардлагатай. Иймд Facebook чатботын агуулгыг системийн мэдээлэлтэй тулгах ажил цаашид хийгдэнэ.

\section{Мэдээлэл цуглуулсан арга}
\label{sec:data-sources}

Сургуулийн мэдээллийг хоёр эх сурвалжаас цуглуулсан.

\textbf{Сургуулийн вэб сайт.} Элсэлт, сургалтын төлбөр, хөтөлбөр зэрэг нийтэд зориулсан мэдээллийг сургуулийн вэб сайтаас авсан. Нийтэд нээлттэй хуудас бүрийг хаяг, огноотой нь хадгалж, хайлтад зориулан богино хэсгүүдэд хуваасан (цаашид «хайлтын хэсэг»).

\textbf{Сургуулийн зүгээс өгсөн мэдээлэл.} Өрөө, анги бүлэг, хичээлийн хуваарь, багш, ажилтны мэдээллийг сургуулиас авсан. Үүний тулд хэрэгтэй мэдээллийг сэдвээр нь жагсааж сургуульд хүргүүлсэн. Энэ мэдээллийг өрөө, багш ба ажилтан, анги бүлэг, хичээл гэсэн дөрвөн хүснэгтэд оруулсан (цаашид «сургуулийн хүснэгт»; \fref{fig:datamodel}). Ресепшний ажилтан эцэг эх, зочдод багшийн сургуулийн имэйл хаягийг өгдөг тул имэйл хаягийг багш, ажилтны хүснэгтэд оруулсан. Хувийн утасны дугаар, хувийн имэйл хаягийг системд оруулаагүй.
% ӨӨРИЙН ҮГЭЭР: хэнээс (сургалтын алба, карьерын төв, олон нийтийн алба, ресепшн) юу авснаа
% нэг өгүүлбэрээр нэм. Facebook чатботыг хариуцдаг багштай уулзсаны дараа мөн нэм.

\begin{figure}[ht]
\centering
\adjustbox{max width=\textwidth}{% Ресепшний хүснэгтийн өгөгдлийн бүтэц — UML классын диаграмм (TikZ).
% Эх: reception_tables.py-ийн Room, Service, Staff, ClassGroup, Lesson класс (2026.11.12;
% Staff.email-ийг ресепшний хэрэгцээгээр нэмсэн). Талбарын нэрийг Монголоор бичив.
% 2026.11.23 (багш: «граф дээрх тоонууд хэрэг байна уу, доор нь тайлбар бий»): холбоосын олонлогийн тоог
% (1, 0..1, 0..*) хасав; холбоо бүрийн утгыг шошго ба 3.3-ын бичвэр илэрхийлнэ.
\begin{tikzpicture}[
    font=\footnotesize,
    cls/.style={draw=nmGreenD, line width=0.7pt, rectangle split, rectangle split parts=2, align=left,
                text width=3.6cm, inner sep=3pt, rectangle split part fill={nmGreenS, nmGreen!35}},
    rel/.style={semithick},
    lab/.style={font=\scriptsize, inner sep=1.5pt}]

\node[cls] (room) at (0,0) {\centering\textbf{Өрөө}
  \nodepart{second} дугаар\\давхар\\зориулалт\\тэнхим};
\node[cls] (service) at (7.2,0) {\centering\textbf{Алба, үйлчилгээ}
  \nodepart{second} нэр\\ажлын цаг\\албаны холбоо барих};
\node[cls] (staff) at (-7.2,0) {\centering\textbf{Багш, ажилтан}
  \nodepart{second} нэр\\нэгж, тэнхим\\албан тушаал\\уулзах цаг\\сургуулийн имэйл};
\node[cls] (group) at (-7.2,-5.2) {\centering\textbf{Анги бүлэг}
  \nodepart{second} бүлгийн нэр\\курс\\тэнхим};
\node[cls] (lesson) at (0,-5.2) {\centering\textbf{Хичээл (хуваарийн мөр)}
  \nodepart{second} гараг\\эхлэх, дуусах цаг\\хичээлийн нэр};

% --- холбоос
\draw[rel] (room.east) -- node[lab, below] {байрлуулна} (service.west);

\draw[rel] (staff.east) -- node[lab, below] {ажлын өрөө} (room.west);

\draw[rel] (group.north) -- node[lab, right] {ангийн багш} (staff.south);

\draw[rel] (group.east) -- node[lab, below] {суралцана} (lesson.west);

\draw[rel] (lesson.north) -- node[lab, right] {орно} (room.south);

\draw[rel] (lesson.north west) -- node[lab, sloped, above] {заана} (staff.south east);

\draw[rel] (group.south) -- ++(0,-0.6) -| node[lab, pos=0.25, below] {үндсэн өрөө} ([xshift=3.4cm]lesson.south) |- ([yshift=-3mm]room.east);
\end{tikzpicture}
}
\caption{Сургуулийн хүснэгтийн өгөгдлийн бүтэц}
\label{fig:datamodel}
\end{figure}

Хоёр эх сурвалжаас цуглуулсан мэдээллийн хэмжээг \tref{tab:collected_data}-д харуулав. Мэдээллийг шинэчлэхдээ шинэ хувилбарыг бэлтгэж, шалгасны дараа нийтэлдэг. Өмнөх хувилбар хадгалагдах тул алдаа гарвал буцааж болно. Цаашид үүнийг сургуулийн хариуцсан ажилтан хийх боломжтой байхаар төлөвлөсөн (\fref{fig:usecase}).

\section{Мэдлэгийн модонд суурилсан хайлт ба хариулт}
\label{sec:rag-method}

Систем хайлтаар баяжуулсан хариултын аргаар (\ref{sec:lit-rag}-р хэсэг) хариулна. Түүний суурь нь цуглуулсан мэдээллийг сэдвээр нь шатлан ангилсан мэдлэгийн мод юм. Энэ хэсэгт модны бүтэц болон модыг хайлт, хариултад хэрхэн ашиглахыг тодорхойлов.

\subsection{Мэдлэгийн модны бүтэц}

Мод гурван түвшинтэй: үндэс нь сургууль, хоёр дахь түвшин нь 11 хүрээ, гурав дахь түвшин нь 31 дэд сэдэв (\fref{fig:tree}, \tref{tab:tree}). Жишээ нь «Элсэлт» хүрээ элсэх шаардлага, шалгалт, бүрдүүлэх материал, төлбөр, бүртгэл, олимпиад гэсэн зургаан дэд сэдэвтэй. Модыг зангилаа бүрийн нэр, дээд зангилааг заасан JSON файл хэлбэрээр хадгалсан.

\begin{figure}[ht]
\centering
\adjustbox{max width=\textwidth}{% Мэдлэгийн модны бүтэц (TikZ). Эх: config/knowledge-topics.json (2026.11.16-нд тулгасан):
% үндэс -> 11 хүрээ -> 31 дэд сэдэв. Зурагт «Элсэлт», «Байршил, өрөө» хоёр хүрээний дэд сэдвийг жишээ болгон задалсан.
% 2026.11.20 (багш нар 11/17: «өнгө бүгд ижил»): хүрээг мэдээллийн эх сурвалжаар нь өнгөөр ялгасан —
% хөх = вэб сайтаас авсан 8 хүрээ, ногоон = сургуулийн зүгээс өгсөн мэдээллийн 3 хүрээ (Хүснэгт 3.1-тэй ижил);
% дэд сэдэв нь хүрээнийхээ өнгөтэй. Доод талд тайлбартай.
% 2026.11.23 (багш: «жишээ болгон задалсан хүрээ өөр өнгөтэй байх хэрэг байна уу, бусад шиг болго»):
% задалсан хоёр хүрээний тод дүүргэлтийг авч бусадтай ижил болгов; гарчгийн «(11)» тоог хасав (тайлбар доор нь).
\begin{tikzpicture}[
    font=\footnotesize,
    node distance=0pt,
    box/.style={draw=nmBlueD, fill=nmBlue, line width=0.7pt, rounded corners=1.5pt, align=left, inner sep=3pt, minimum height=5.2mm},
    dom/.style={box, text width=4.3cm, anchor=west},
    school/.style={draw=nmGreenD, fill=nmGreen},
    subt/.style={box, text width=6.6cm, anchor=west},
    key/.style={minimum width=4mm, minimum height=3mm, inner sep=0pt, line width=0.7pt},
    note/.style={font=\scriptsize, inner sep=1.5pt},
    line/.style={semithick, black!70}]

\node[box, draw=nmVioletD, fill=nmViolet, align=center, text width=2.3cm] (root) at (0,-3.25) {Сургуулийн\\мэдээлэл};

\node[dom]       (d0)  at (2.9,0)     {Сургуулийн тухай};
\node[dom]       (d1)  at (2.9,-0.65) {Элсэлт};
\node[dom]       (d2)  at (2.9,-1.30) {Сургалт, тэнхимүүд};
\node[dom]       (d3)  at (2.9,-1.95) {Оюутны амьдрал};
\node[dom]       (d4)  at (2.9,-2.60) {Карьер, хамтын ажиллагаа};
\node[dom]       (d5)  at (2.9,-3.25) {Мэдээ};
\node[dom]       (d6)  at (2.9,-3.90) {Арга хэмжээ};
\node[dom]       (d7)  at (2.9,-4.55) {Холбоо барих};
\node[dom, school] (d8)  at (2.9,-5.20) {Байршил, өрөө};
\node[dom, school] (d9)  at (2.9,-5.85) {Багш, ажилтан};
\node[dom, school] (d10) at (2.9,-6.50) {Анги, хичээлийн хуваарь};

\foreach \n in {d0,d1,d2,d3,d4,d6,d7,d8,d9,d10}
    \draw[line] (root.east) -- ++(0.45,0) |- (\n.west);
\draw[line] (root.east) -- (d5.west);

\node[subt] (a1) at (8.6,0.35)  {Элсэх шаардлага};
\node[subt] (a2) at (8.6,-0.30) {Шалгалт, эссэ, ярилцлага};
\node[subt] (a3) at (8.6,-0.95) {Бүрдүүлэх материал};
\node[subt] (a4) at (8.6,-1.60) {Төлбөр, хураамж};
\node[subt] (a5) at (8.6,-2.25) {Бүртгэл, элсэлтийн журам};
\node[subt] (a6) at (8.6,-2.90) {Олимпиад, холбогдох элсэлтийн нөхцөл};
\foreach \n in {a1,a2,a3,a4,a5,a6}
    \draw[line] (d1.east) -- ++(0.45,0) |- (\n.west);

\node[subt, school] (l1) at (8.6,-4.55) {Давхар бүрийн өрөө};
\node[subt, school] (l2) at (8.6,-5.20) {Ресепшн, үйлчилгээний өрөө};
\node[subt, school] (l3) at (8.6,-5.85) {Тэнхимийн өрөө};
\foreach \n in {l1,l3}
    \draw[line] (d8.east) -- ++(0.45,0) |- (\n.west);
\draw[line] (d8.east) -- (l2.west);

\node[font=\scriptsize\bfseries, anchor=south] at (root.north |- 0,0.75) {Үндэс};
\node[font=\scriptsize\bfseries, anchor=south west] at (2.9,0.75) {Хүрээ};
\node[font=\scriptsize\bfseries, anchor=south west] at (8.6,0.75) {Дэд сэдэв (жишээ)};

% --- өнгөний тайлбар
\node[key, draw=nmBlueD, fill=nmBlue] (k1) at (3.1,-7.3) {};
\node[note, anchor=west] at (k1.east) {вэб сайтаас авсан мэдээлэл};
\node[key, draw=nmGreenD, fill=nmGreen] (k2) at (8.8,-7.3) {};
\node[note, anchor=west] at (k2.east) {сургуулийн зүгээс өгсөн мэдээлэл};
\end{tikzpicture}
}
\caption{Мэдлэгийн модны бүтэц: 11 хүрээ ба хоёр хүрээний дэд сэдэв}
\label{fig:tree}
\end{figure}

Цуглуулсан мэдээлэл бүрийг модны зангилаатай холбосон. Вэб сайтын хуудас бүрийг хаягаар нь нэг зангилаанд оноосон тул хайлтын хэсэг бүр аль хүрээ, аль дэд сэдэвт хамаарах нь тодорхой байна. Сургуулийн хүснэгтийг байршил, багш, анги гэсэн гурван хүрээнд холбосон. Мэдээ, арга хэмжээ, холбоо барих хүрээг дэд сэдэвт хуваагаагүй.

Модыг сургуулийн вэб сайтын цэс болон сургуулийн зүгээс өгсөн мэдээлэлд тулгуурлан гараар байгуулсан. Энэ нь сургуулийн албан ёсны ангилал биш, энэ ажлын хүрээнд боловсруулсан ангилал юм. Мод нь хайлтаас гадна аль дэд сэдвийн мэдээллийг нэмж цуглуулах шаардлагатайг харуулдаг. Сургуульд хүргүүлсэн хэрэгтэй мэдээллийн жагсаалтыг үүний дагуу гаргасан.

\begin{table}[H]
\centering
\caption{Мэдлэгийн модны 11 хүрээ (2026.11.15-ны байдлаар)}
\label{tab:tree}
\begin{tabular}{|p{5.6cm}|r|p{5.6cm}|}
\hline
\textbf{Хүрээ} & \textbf{Дэд сэдэв} & \textbf{Мэдээллийн эх сурвалж} \\
\hline
Сургуулийн тухай & 3 & Вэб сайт \\
\hline
Элсэлт & 6 & Вэб сайт \\
\hline
Сургалт, тэнхимүүд & 7 & Вэб сайт \\
\hline
Оюутны амьдрал & 4 & Вэб сайт \\
\hline
Карьер, хамтын ажиллагаа & 3 & Вэб сайт \\
\hline
Мэдээ & 0 & Вэб сайт \\
\hline
Арга хэмжээ & 0 & Вэб сайт \\
\hline
Холбоо барих & 0 & Вэб сайт \\
\hline
Байршил, өрөө & 3 & Сургуулийн зүгээс өгсөн мэдээлэл \\
\hline
Багш, ажилтан & 3 & Сургуулийн зүгээс өгсөн мэдээлэл \\
\hline
Анги, хичээлийн хуваарь & 2 & Сургуулийн зүгээс өгсөн мэдээлэл \\
\hline
\textbf{Нийт} & \textbf{31} & \\
\hline
\end{tabular}
\end{table}

\subsection{Хайлт ба хариулт}

Асуултад хариулахдаа мэдлэгийн модыг дараах таван алхмаар ашиглана.
\begin{enumerate}
    \item \textbf{Хүрээг тодорхойлох.} Асуулт 11 хүрээний алинд хамаарахыг хүрээ бүрийн түлхүүр үгээр тодорхойлно. Жишээ нь «элсэлтийн төлбөр хэд вэ» гэсэн асуулт «Элсэлт» хүрээнд хамаарна. Хоёр ба түүнээс олон хүрээ тэнцүү таарвал асуултыг өргөн хүрээний асуулт гэж үзэж, хүрээг тодорхойлохгүй.
    \item \textbf{Хайх.} Бүх хайлтын хэсгээс хоёр аргаар зэрэг хайна. BM25 нь асуултын үгтэй таарах хэсгийг \cite{robertson2009bm25}, E5 загвар нь утгаараа ойр хэсгийг олно \cite{wang2024me5}. Хоёр жагсаалтыг \acs{RRF} аргаар нэгтгэнэ \cite{cormack2009rrf}. Энэ арга хэсэг бүрд жагсаалт тус бүрд эзэлсэн байраар нь оноо өгдөг.
    \item \textbf{Хүрээний эх сурвалжид нэмэлт оноо өгөх.} Нэгтгэсэн жагсаалтад асуултын хүрээнд хамаарах хэсгүүдэд нэмэлт оноо өгнө. Бусад хэсгийг хасахгүй тул хүрээг буруу тодорхойлсон үед ч зөв эх сурвалж жагсаалтад үлдэнэ (\ref{sec:lit-tree}-р хэсэг).
    \item \textbf{Сонгож шалгах.} Жагсаалтаас хариулт болох хэсгийг дүрмээр сонгож, асуултын он, нөхцөл, шаардлагатай талбарыг агуулж байгаа эсэхийг шалгана. Энэ алхамд хэлний загвар оролцохгүй.
    \item \textbf{Хариулах эсвэл чиглүүлэх.} Шаардлага хангасан бол эх сурвалжийн өгүүлбэрийг холбоосын хамт үгчлэн харуулна. Хангаагүй бол асуултын хүрээгээр хэнд, хаана хандахыг заана.
\end{enumerate}

Хүрээгээр нэмэлт оноо өгөх хэсгийг (нэг ба гуравдугаар алхам) хэрэгжүүлсэн бөгөөд түүнийг идэвхжүүлэх эсэхийг системийн үнэлгээний дүнгээр шийднэ. Тавдугаар алхамд систем одоогоор ресепшний ажилтан ба сургуулийн ерөнхий утсыг заадаг. Сэдэв бүрийг хариуцах алба, өрөөний жагсаалтыг ресепшнээс авч сургуулиар батлуулсны дараа систем тодорхой алба, өрөөг нэрлэнэ.

\section{Системийн ерөнхий бүтэц}
\label{sec:system}

Системийг вэб програм хэлбэрээр хэрэгжүүлсэн. Иймд сургуулийн нэгдүгээр давхрын ресепшний дэлгэц дээр ч, гар утасны хөтчөөр ч нэг ижил систем нээгдэнэ. Гар утаснаас нээгддэг байхыг сургалтын алба санал болгосон (\ref{sec:intro-background}-р хэсэг). Систем одоогоор хөгжүүлэлтийн компьютер дээр ажиллаж байгаа бөгөөд эхний туршилтыг ресепшний дэлгэц дээр хийнэ.

Системийн ерөнхий бүтцийг \fref{fig:system}-т харуулав. Систем асуултыг таван алхмаар боловсруулж хариулт гаргана.

\begin{figure}[ht]
\centering
\adjustbox{max width=\textwidth}{% Системийн бүтэц — асуултаас хариулт хүртэл (TikZ).
% Эх: локал системийн /data#tree хуудасны зураг (docs/diagrams/Системийн_бүтэц.pdf, 11/06),
% 11/09-нөөс асаасан ресепшний хүснэгт ба «хэнд, хаана хандах» давхаргыг нэмсэн.
% Байрлалыг (x,y) координатаар өгсөн: нэг хайрцгийг зөөхөд тоог нь л өөрчилнө.
% 2026.11.16 (багшийн зөвлөгөө): өнгө нэмсэн — хөх = системийн алхам, улбар шар = гадны
% үйлчилгээ, ногоон = мэдээлэл, ягаан = хэнд, хаана хандахыг заах. Доод талд тайлбартай.
% 2026.11.20 (багш нар 11/17: «өнгө бүгд ижил»): өнгийг Figures/palette.tex-ийн тод палитраар сольсон.
% 2026.11.23 (багш: «граф дээрх тоонууд хэрэг байна уу, доор нь тайлбар бий»): баганын гарчгийн 1–5 дугаар, «Мэдлэгийн мод»-ны «11 хүрээ, 31 дэд сэдэв» тоог хасав (3.4, Хүснэгт 3.1-д бий).
\begin{tikzpicture}[
    font=\footnotesize,
    box/.style={draw=nmBlueD, fill=nmBlue, line width=0.7pt, rounded corners=2pt, align=center, text width=2.3cm,
                minimum height=0.95cm, inner sep=2pt},
    ext/.style={box, draw=nmOrangeD, fill=nmOrange},
    data/.style={box, draw=nmGreenD, fill=nmGreen},
    ref/.style={box, draw=nmVioletD, fill=nmViolet},
    key/.style={minimum width=4mm, minimum height=3mm, inner sep=0pt, line width=0.7pt},
    head/.style={font=\footnotesize\bfseries, align=center},
    arr/.style={-{Stealth[length=2mm]}, semithick},
    note/.style={font=\scriptsize, inner sep=1.5pt}]

% --- баганын гарчиг
\node[head] at (0,1.0)     {Асуух};
\node[head] at (3.1,1.0)   {Ангилах};
\node[head] at (6.2,1.0)   {Хайх};
\node[head] at (9.3,1.0)   {Сонгох};
\node[head] at (13.0,1.0)  {Шалгаж хариулах};

% --- 1. Асуух
\node[box]  (voice) at (0,0)      {Дуугаар асуух};
\node[ext]  (stt)   at (0,-1.7)   {Яриа таних\\{\scriptsize Gemini 3.5 Transcribe}};
\node[box]  (text)  at (0,-3.4)   {Асуултын текст};
\node[box]  (typed) at (0,-5.1)   {Бичгээр асуух};

% --- 2. Чиглүүлэх
\node[box]  (clar)  at (3.1,-1.7) {Тодруулж асуух};
\node[box]  (route) at (3.1,-3.4) {Асуулт ангилагч};

% --- 3. Хайх
\node[box]  (bm25)  at (6.2,-1.7)  {BM25\\{\scriptsize үгээр хайх}};
\node[box]  (e5)    at (6.2,-3.4)  {E5\\{\scriptsize утгаар хайх}};
\node[box]  (table) at (6.2,-5.1)  {Сургуулийн хүснэгт\\{\scriptsize өрөө, анги, багш, хуваарь}};

% --- 4. Сонгох
\node[box]  (rrf)   at (9.3,-1.7) {RRF\\{\scriptsize хоёр дүнг нэгтгэх}};
\node[box]  (sel)   at (9.3,-3.4) {Хэсэг сонгох дүрэм};

% --- 5. Шалгаж хариулах
\node[ref]  (refer) at (13.0,-1.7) {Хэнд, хаана\\хандахыг заах};
\node[box]  (guard) at (13.0,-3.4) {Хариултыг шалгах};
\node[box]  (ans)   at (13.0,-5.1) {Эх сурвалжтай\\хариулт};
\node[ext]  (tts)   at (13.0,-7.05) {Дуугаар унших\\{\scriptsize Edge TTS}};

% --- сумнууд
\draw[arr] (voice) -- (stt);
\draw[arr] (stt) -- (text);
\draw[arr] (typed) -- (text);
\draw[arr] (text) -- (route);
\draw[arr] (route) -- node[note, left] {тодорхойгүй} (clar);
\draw[arr] (route.east) -- ++(0.25,0) |- (bm25.west);
\draw[arr] (route) -- (e5);
\draw[arr] (route.east) -- ++(0.25,0) |- (table.west);
\draw[arr] (bm25) -- (rrf);
\draw[arr] (e5.east) -- ++(0.25,0) |- ([yshift=-2.5mm]rrf.west);
\draw[arr] (rrf) -- (sel);
\draw[arr] (rrf) -- node[note, above] {олдоогүй} (refer);
\draw[arr] (sel) -- (guard);
\draw[arr] (table) -- node[note, above, pos=0.3] {мөр олдвол} (ans);
\draw[arr] (table.south) -- (6.2,-6.2) -- node[note, above, pos=0.4] {мөр олдоогүй} (14.7,-6.2) -- (14.7,-1.7) -- (refer.east);
\draw[arr] (guard) -- node[note, right] {хангаагүй} (refer);
\draw[arr] (guard) -- node[note, right] {хангасан} (ans);
\draw[arr] (ans) -- (tts);

% --- Мэдээлэл бэлтгэх (асуултаас өмнө тусдаа хийгдэнэ)
\draw[dashed] (-1.35,-7.95) -- (14.75,-7.95);
\node[head, anchor=west] at (-1.35,-8.3) {Мэдээлэл бэлтгэх};
\node[data] (web)    at (0,-9.2)     {Сургуулийн\\вэб сайт};
\node[data] (office) at (0,-10.9)    {Сургуулийн өгсөн жагсаалт, хуваарь};
\node[data] (store)  at (4.2,-10.05) {Баримтын сан\\{\scriptsize огноотой хувилбар}};
\node[data] (chunk)  at (8.4,-9.2)   {Хайлтын хэсгүүд};
\node[data] (rtab)   at (8.4,-10.9)  {Сургуулийн хүснэгт};
\node[data] (tree)   at (13.0,-10.05) {Мэдлэгийн мод};

\draw[arr] (web.east) -- ++(0.6,0) |- ([yshift=2mm]store.west);
\draw[arr] (office.east) -- ++(0.6,0) |- ([yshift=-2mm]store.west);
\draw[arr] (store.east) -- ++(0.6,0) |- (chunk.west);
\draw[arr] (store.east) -- ++(0.6,0) |- (rtab.west);
\draw[arr] (chunk.east) -- ++(0.9,0) |- ([yshift=2mm]tree.west);
\draw[arr] (rtab.east) -- ++(0.9,0) |- ([yshift=-2mm]tree.west);

% --- өнгөний тайлбар
\node[key, draw=nmBlueD, fill=nmBlue] (k1) at (0.6,-12.2) {};
\node[note, anchor=west] at (k1.east) {системийн алхам};
\node[key, draw=nmOrangeD, fill=nmOrange] (k2) at (4.2,-12.2) {};
\node[note, anchor=west] at (k2.east) {гадны үйлчилгээ};
\node[key, draw=nmGreenD, fill=nmGreen] (k3) at (7.8,-12.2) {};
\node[note, anchor=west] at (k3.east) {мэдээлэл};
\node[key, draw=nmVioletD, fill=nmViolet] (k4) at (10.4,-12.2) {};
\node[note, anchor=west] at (k4.east) {хэнд, хаана хандахыг заах};
\end{tikzpicture}
}
\caption{Системийн бүтэц: асуултаас хариулт хүртэл}
\label{fig:system}
\end{figure}

\begin{enumerate}
    \item \textbf{Асуух.} Хэрэглэгч асуултаа бичнэ эсвэл микрофоноор хэлнэ. Хэлсэн асуултыг яриа таних үйлчилгээ текст болгоно. Яриа гадны үйлчилгээ рүү очихыг дэлгэц дээр мэдэгдэнэ. Хэрэглэгч системийн таньсан текстийг илгээхээс өмнө засаж болно. Яриа таних үйлчилгээ ажиллахгүй үед асуултаа бичиж асууна.
    \item \textbf{Ангилах.} Систем асуултын төрлийг тодорхойлно. Өрөө, анги, багш, хуваарийн асуултад сургуулийн хүснэгтээс, бусад асуултад хайлтаар хариулна. Асуулт тодорхойгүй бол хэрэглэгчээс тодруулж асууна.
    \item \textbf{Хайх.} Мэдлэгийн модонд суурилсан хайлтаар (\ref{sec:rag-method}-р хэсэг) асуултад тохирох хайлтын хэсгүүдийг олно.
    \item \textbf{Сонгох.} Олдсон хэсгүүдээс хариулт болох хэсгийг дүрмээр сонгоно.
    \item \textbf{Шалгаж хариулах.} Сонгосон хэсэг шаардлага хангасан бол эх сурвалжийн өгүүлбэрийг холбоосын хамт харуулж, хэрэглэгч хүсвэл дуугаар уншиж өгнө. Хангаагүй эсвэл тохирох эх сурвалж, хүснэгтийн мөр олдоогүй бол хэнд, хаана хандахыг заана.
\end{enumerate}

Хариултыг дуугаар уншихад Microsoft Edge-ийн онлайн дуу үүсгэх үйлчилгээг edge-tts сангаар ашигласан \cite{edge_tts}. Яриа таних ба дуу үүсгэх нь гадны үйлчилгээ бөгөөд бусад алхам системийн сервер дээр ажиллана.

Нэг асуултад хариулах үйл явцыг \fref{fig:activity}-т үйл ажиллагааны диаграммаар харуулав.

\begin{figure}[p]
\centering
\adjustbox{max width=\textwidth, max totalheight=0.88\textheight}{% Асуултад хариулах үйл явц — UML идэвхийн диаграмм (TikZ).
% Эх: docs/diagrams/UML_идэвхийн_диаграмм.pdf (11/06)-ийг хялбарчилсан; ресепшний хүснэгт
% ба «хэнд, хаана хандах» алхмыг нэмсэн. Байрлалыг (x,y) координатаар өгсөн.
% 2026.11.15 (3): 5 хянагчийн шалгалтын D1–D4 — таньсан текстийг засах, тодруулж асуух,
% хүснэгтийн мөр олдоогүй үед чиглүүлэх, дуугаар сонсох эсэхийг сонгох алхмыг нэмсэн.
\begin{tikzpicture}[
    font=\footnotesize,
    act/.style={draw=nmBlueD, fill=nmBlue, line width=0.7pt, rounded corners=6pt, align=center, text width=4.2cm,
                minimum height=0.85cm, inner sep=2pt},
    ext/.style={act, draw=nmOrangeD, fill=nmOrange},
    ref/.style={act, draw=nmVioletD, fill=nmViolet},
    dec/.style={draw, fill=white, diamond, aspect=2, minimum width=1.1cm, minimum height=0.55cm, inner sep=0pt},
    arr/.style={-{Stealth[length=2mm]}, semithick},
    q/.style={font=\scriptsize\itshape},
    g/.style={font=\scriptsize, inner sep=1.5pt}]

\node[circle, fill=black, minimum size=3.2mm, inner sep=0pt] (start) at (0,0) {};
\node[act] (ask)   at (0,-1.0)  {Асуултаа ярих эсвэл бичих};
\node[dec] (dvoice) at (0,-2.3) {};
\node[q, anchor=east] at (-0.7,-2.3) {Дуугаар асуусан уу?};
\node[ext] (stt)   at (5.2,-2.3) {Яриаг текст болгох\\{\scriptsize Gemini 3.5 Transcribe}};
\node[act] (edit)  at (5.2,-3.7) {Таньсан текстийг засах};
\node[act] (route) at (0,-3.7)  {Асуултыг ангилах};
\node[act] (clar)  at (-5.2,-3.7) {Тодруулж асуух};
\node[dec] (dtab)  at (0,-5.0)  {};
\node[q, anchor=west] at (0.7,-5.0) {Өрөө, анги, багш, хуваарийн асуулт уу?};
\node[act] (table) at (-5.2,-6.4) {Сургуулийн хүснэгтээс мөрийг хайх};
\node[act] (search) at (0,-6.4) {BM25 ба E5-аар хайж,\\RRF-ээр нэгтгэх};
\node[act] (select) at (0,-7.8) {Дүрмээр хэсэг сонгох};
\node[act] (guard) at (0,-9.2)  {Сонгосон хэсгийг шалгах};
\node[dec] (dpass) at (0,-10.5) {};
\node[q, anchor=east] at (-0.5,-10.1) {Хариулт олдсон уу?};
\node[ref] (refer) at (5.2,-10.5) {Хэнд, хаана хандахыг заах};
\node[act] (answer) at (0,-11.9) {Эх сурвалжийн үгээр\\хариулт бүрдүүлэх};
\node[act] (show)  at (0,-13.3) {Хариулт, эх сурвалжийг харуулах};
\node[dec] (dlisten) at (0,-14.6) {};
\node[q, anchor=east] at (-0.7,-14.6) {Дуугаар сонсох уу?};
\node[ext] (tts)   at (0,-15.9) {Хариултыг дуугаар унших\\{\scriptsize Edge TTS}};
\node[circle, draw, minimum size=4.4mm, inner sep=0pt] (end) at (0,-17.1) {};
\node[circle, fill=black, minimum size=2.6mm, inner sep=0pt] at (end) {};

\draw[arr] (start) -- (ask);
\draw[arr] (ask) -- (dvoice);
\draw[arr] (dvoice) -- node[g, above] {[тийм]} (stt);
\draw[arr] (dvoice) -- node[g, left] {[үгүй]} (route);
\draw[arr] (stt) -- (edit);
\draw[arr] (edit) -- (route);
\draw[arr] (route) -- (clar);
\node[g] at (-2.6,-3.1) {[тодорхойгүй]};
\draw[arr] (clar.north) |- (ask.west);
\draw[arr] (route) -- (dtab);
\draw[arr] (dtab.west) -| node[g, above, pos=0.25] {[тийм]} (table.north);
\draw[arr] (dtab) -- node[g, right] {[үгүй]} (search);
\draw[arr] (search) -- (select);
\draw[arr] (select) -- (guard);
\draw[arr] (table.south) |- (dpass.west);
\draw[arr] (guard) -- (dpass);
\draw[arr] (dpass) -- node[g, above] {[үгүй]} (refer);
\draw[arr] (dpass) -- node[g, right] {[тийм]} (answer);
\draw[arr] (answer) -- (show);
\draw[arr] (refer.south) |- (show.east);
\draw[arr] (show) -- (dlisten);
\draw[arr] (dlisten) -- node[g, right] {[тийм]} (tts);
\draw[arr] (dlisten.east) -- ++(3.0,0) node[g, above, pos=0.5] {[үгүй]} |- (end.east);
\draw[arr] (tts) -- (end);
\end{tikzpicture}
}
\caption{Асуултад хариулах үйл явцын үйл ажиллагааны диаграмм}
\label{fig:activity}
\end{figure}

Дуугаар асуусан нэг асуултад системийн хэсгүүд ямар дарааллаар мэдээлэл солилцохыг \fref{fig:sequence}-д харуулав.

\begin{figure}[p]
\centering
\adjustbox{max width=\textwidth, max totalheight=0.88\textheight}{% Дуугаар асуусан нэг асуултын дарааллын (sequence) диаграмм (TikZ).
% Эх: demo_server.py — /api/voice/transcribe, /api/route, /api/voice/speak;
% reception_tables.py (хүснэгтийн зам), website_retrieval.py (хайлт), referral.py.
% Босоо шугам = оролцогч, хэвтээ сум = мессеж; y тэнхлэг доошоо цаг.
\begin{tikzpicture}[
    font=\footnotesize,
    obj/.style={draw=nmBlueD, fill=nmBlue, line width=0.7pt, align=center, text width=2.2cm, minimum height=0.9cm},
    ext/.style={obj, draw=nmOrangeD, fill=nmOrange},
    usr/.style={obj, draw=nmTealD, fill=nmTeal},
    dat/.style={obj, draw=nmGreenD, fill=nmGreen},
    life/.style={dashed, black!60},
    msg/.style={-{Stealth[length=2mm]}, semithick},
    ret/.style={-{Stealth[length=2mm]}, semithick, dashed},
    lab/.style={font=\scriptsize, inner sep=1.5pt, above},
    frame/.style={draw, black!70},
    tag/.style={font=\scriptsize\bfseries, anchor=north west, inner sep=2pt, draw, black!70, fill=white}]

% --- оролцогчид (x байрлал)
\def\xu{0}\def\xb{2.9}\def\xs{5.8}\def\xg{8.7}\def\xk{11.6}\def\xt{14.5}
\node[usr] (u) at (\xu,0) {Хэрэглэгч};
\node[obj] (b) at (\xb,0) {Вэб хуудас};
\node[obj] (s) at (\xs,0) {Сервер\\{\scriptsize ангилагч}};
\node[ext] (g) at (\xg,0) {Яриа таних\\{\scriptsize Gemini}};
\node[dat] (k) at (\xk,0) {Сургуулийн хүснэгт,\\баримтын сан};
\node[ext, text width=2.6cm] (t) at (\xt,0) {Текстийг дуу болгох\\{\scriptsize Edge TTS}};
\foreach \n/\x in {u/\xu,b/\xb,s/\xs,g/\xg,k/\xk,t/\xt} \draw[life] (\n.south) -- (\x,-14.3);

% --- 1. яриа таних
\draw[msg] (\xu,-1.1) -- node[lab] {асуултаа хэлнэ} (\xb,-1.1);
\draw[msg] (\xb,-1.7) -- node[lab] {бичлэг} (\xs,-1.7);
\draw[msg] (\xs,-2.3) -- node[lab] {бичлэг, нэрийн жагсаалт} (\xg,-2.3);
\draw[ret] (\xg,-2.9) -- node[lab] {текст} (\xs,-2.9);
\draw[ret] (\xs,-3.5) -- node[lab] {текст} (\xb,-3.5);
\node[font=\scriptsize\itshape, anchor=west] at (\xg+0.12,-3.3) {хязгаарт хүрвэл: бичгээр асууна};
\draw[msg] (\xu,-4.1) -- node[lab] {засаж, илгээнэ} (\xb,-4.1);
\draw[msg] (\xb,-4.7) -- node[lab] {асуултын текст} (\xs,-4.7);

% --- 2. хариулт хайх (alt)
\draw[frame] (\xs-1.3,-5.2) rectangle (\xk+1.4,-10.4);
\node[tag] at (\xs-1.3,-5.2) {alt};
\node[font=\scriptsize\itshape, anchor=west] at (\xs-0.55,-5.45) {[өрөө, анги, багш, хуваарийн асуулт]};
\draw[msg] (\xs,-6.2) -- node[lab] {хүснэгтээс харах} (\xk,-6.2);
\draw[ret] (\xk,-6.8) -- node[lab] {мөр} (\xs,-6.8);
\draw[frame, dashed] (\xs-1.3,-7.2) -- (\xk+1.4,-7.2);
\node[font=\scriptsize\itshape, anchor=west] at (\xs-0.55,-7.45) {[бусад асуулт]};
\draw[msg] (\xs,-8.2) -- node[lab] {BM25 + E5 хайлт} (\xk,-8.2);
\draw[ret] (\xk,-8.8) -- node[lab] {олдсон хэсгүүд} (\xs,-8.8);
\draw[msg] (\xs,-9.4) -- ++(0.9,0) -- ++(0,-0.5) -- (\xs,-9.9);
\node[font=\scriptsize, anchor=west] at (\xs+0.95,-9.65) {хэсэг сонгож, шалгана};

% --- 3. хариулт
\draw[ret] (\xs,-10.9) -- node[lab] {хариулт + эх сурвалж} (\xb,-10.9);
\node[font=\scriptsize\itshape, anchor=west] at (\xs+0.1,-11.15) {олдоогүй бол: хэнд, хаана хандах};
\draw[ret] (\xb,-11.5) -- node[lab] {дэлгэцэнд} (\xu,-11.5);
\node[font=\scriptsize\itshape, anchor=west] at (\xb+0.1,-11.95) {хэрэглэгч сонсохыг хүсвэл:};

% --- 4. дуугаар унших
\draw[msg] (\xb,-12.6) -- node[lab] {хариултын текст} (\xs,-12.6);
\draw[msg] (\xs,-13.1) -- node[lab, pos=0.75] {текст} (\xt,-13.1);
\draw[ret] (\xt,-13.6) -- node[lab, pos=0.25] {дуу} (\xs,-13.6);
\draw[ret] (\xs,-14.1) -- node[lab] {дуу} (\xb,-14.1);
\end{tikzpicture}
}
\caption{Дуугаар асуусан асуултын дарааллын диаграмм}
\label{fig:sequence}
\end{figure}

\fullinput{Chapters/full/ch3_evaluation_and_stt_choice}
